% GENERATED FILE. Edit canonical lessons, then run npm run generate:books.
% canonical-source-hash: fnv1a64:2e0611d8ce44b189
% canonical-lessons: HI-C519-sanduq, HI-C519-makri, HI-C519-raddi, HI-C519-buharna, HI-C519-majna

\chapter{Clearing the Clutter}
\label{ch:hi-clearing-the-clutter}
\hlchaptermodality{519}

\begin{chapteropening}
\textbf{By the end of this chapter:} \emph{I can talk about dust, old papers and scrubbing.}

Complete the last lesson of chapter 519: I can talk about dust, old papers and scrubbing.
\end{chapteropening}

\section[\texorpdfstring{\hi{संदूक़} (sandūq)}{संदूक़ (sandūq)}]{\hi{संदूक़} (sandūq) — a trunk; a chest}

\label{lesson:HI-C519-sanduq}



\begin{quote}
Before the new one: say the Hindi for a peg, then the Hindi for to hang up.
\end{quote}

\subsection*{You'll want to know: \hi{संदूक़}}
\textbf{\hi{संदूक़}} — \emph{sandūq} — \textquotedblleft{}a trunk; a chest\textquotedblright{}.

\begin{quote}
\hi{दादी} \hi{के} \hi{पुराने} \hi{कपड़े} \hi{एक} \hi{संदूक़} \hi{में} \hi{हैं।}
\end{quote}

\begin{quote}
\emph{Dādī ke purāne kapṛe ek sandūq meṁ haiṁ.}
\end{quote}

\begin{quote}
Grandmother's old clothes are in a trunk.
\end{quote}

Masculine.

\begin{cousinweb}
\textbf{\hi{संदूक़}} (\emph{sandūq}) is usually traced to Arabic \emph{ṣandūq}, \textquotedblleft{}a chest, a box\textquotedblright{}, through Persian.
\end{cousinweb}

\begin{grammarlens}[title={What you've built}]
One more word for describing what you lived and felt.
\end{grammarlens}

\subsection*{Your turn}
\begin{itemize}
\raggedright
  \item \emph{Say it:} \emph{sandūq}
  \item \emph{Say it:} \emph{sandūq}, once more
  \item \emph{Say it:} \emph{Dādī ke purāne kapṛe ek sandūq meṁ haiṁ.} — Grandmother's old clothes are in a trunk.
  \item \emph{Recall:} say the Hindi for a peg, then the Hindi for to hang up, then say \emph{sandūq} again
\end{itemize}

\begin{tcolorbox}[breakable,colback=teal!4,colframe=teal!35!black,title={Before you move on}]
What does \hi{संदूक़} mean? (\textquotedblleft{}A trunk; a chest\textquotedblright{}.) Say it once more.
\end{tcolorbox}
\section[\texorpdfstring{\hi{मकड़ी} (makṛī)}{मकड़ी (makṛī)}]{\hi{मकड़ी} (makṛī) — a spider}

\label{lesson:HI-C519-makri}



\begin{quote}
Before the new one: say the Hindi for to hang up, then the Hindi for a trunk.
\end{quote}

\subsection*{You'll want to know: \hi{मकड़ी}}
\textbf{\hi{मकड़ी}} — \emph{makṛī} — \textquotedblleft{}a spider\textquotedblright{}.

\begin{quote}
\hi{कोने} \hi{में} \hi{एक} \hi{बड़ी} \hi{मकड़ी} \hi{बैठी} \hi{है।}
\end{quote}

\begin{quote}
\emph{Kone meṁ ek baṛī makṛī baiṭhī hai.}
\end{quote}

\begin{quote}
A big spider is sitting in the corner.
\end{quote}

Feminine.

\begin{cousinweb}
\textbf{\hi{मकड़ी}} (\emph{makṛī}) is a tadbhava, inherited from Sanskrit \emph{markaṭa}, \textquotedblleft{}a spider\textquotedblright{} (CDIAL 9883).
\end{cousinweb}

\begin{grammarlens}[title={What you've built}]
One more word for describing what you lived and felt.
\end{grammarlens}

\subsection*{Your turn}
\begin{itemize}
\raggedright
  \item \emph{Say it:} \emph{makṛī}
  \item \emph{Say it:} \emph{makṛī}, once more
  \item \emph{Say it:} \emph{Kone meṁ ek baṛī makṛī baiṭhī hai.} — A big spider is sitting in the corner.
  \item \emph{Recall:} say the Hindi for to hang up, then the Hindi for a trunk, then say \emph{makṛī} again
\end{itemize}

\begin{tcolorbox}[breakable,colback=teal!4,colframe=teal!35!black,title={Before you move on}]
What does \hi{मकड़ी} mean? (\textquotedblleft{}A spider\textquotedblright{}.) Say it once more.
\end{tcolorbox}
\section[\texorpdfstring{\hi{रद्दी} (raddī)}{रद्दी (raddī)}]{\hi{रद्दी} (raddī) — waste paper; scrap}

\label{lesson:HI-C519-raddi}



\begin{quote}
Before the new one: say the Hindi for a trunk, then the Hindi for a spider.
\end{quote}

\subsection*{You'll want to know: \hi{रद्दी}}
\textbf{\hi{रद्दी}} — \emph{raddī} — \textquotedblleft{}waste paper; scrap\textquotedblright{}.

\begin{quote}
\hi{हमने} \hi{पुराने} \hi{अख़बार} \hi{रद्दी} \hi{में} \hi{बेच} \hi{दिए।}
\end{quote}

\begin{quote}
\emph{Hamne purāne akhbār raddī meṁ bec diye.}
\end{quote}

\begin{quote}
We sold the old newspapers as waste paper.
\end{quote}

Feminine.

\begin{cousinweb}
\textbf{\hi{रद्दी}} (\emph{raddī}) is usually traced to Persian \emph{raddī}, \textquotedblleft{}rejected, worthless\textquotedblright{}. The Persian word is built on Arabic \emph{radd}, \textquotedblleft{}rejection\textquotedblright{}.
\end{cousinweb}

\begin{grammarlens}[title={What you've built}]
One more word for describing what you lived and felt.
\end{grammarlens}

\subsection*{Your turn}
\begin{itemize}
\raggedright
  \item \emph{Say it:} \emph{raddī}
  \item \emph{Say it:} \emph{raddī}, once more
  \item \emph{Say it:} \emph{Hamne purāne akhbār raddī meṁ bec diye.} — We sold the old newspapers as waste paper.
  \item \emph{Recall:} say the Hindi for a trunk, then the Hindi for a spider, then say \emph{raddī} again
\end{itemize}

\begin{tcolorbox}[breakable,colback=teal!4,colframe=teal!35!black,title={Before you move on}]
What does \hi{रद्दी} mean? (\textquotedblleft{}Waste paper; scrap\textquotedblright{}.) Say it once more.
\end{tcolorbox}
\section[\texorpdfstring{\hi{बुहारना} (buhārnā)}{बुहारना (buhārnā)}]{\hi{बुहारना} (buhārnā) — to sweep (with a broom)}

\label{lesson:HI-C519-buharna}



\begin{quote}
Before the new one: say the Hindi for a spider, then the Hindi for waste paper.
\end{quote}

\subsection*{You'll want to know: \hi{बुहारना}}
\textbf{\hi{बुहारना}} — \emph{buhārnā} — \textquotedblleft{}to sweep (with a broom)\textquotedblright{}.

\begin{quote}
\hi{दादी} \hi{हर} \hi{सुबह} \hi{आँगन} \hi{बुहारती} \hi{हैं।}
\end{quote}

\begin{quote}
\emph{Dādī har subah āṁgan buhārtī haiṁ.}
\end{quote}

\begin{quote}
Grandmother sweeps the courtyard every morning.
\end{quote}

Past, with \hi{ने} (he or she …): \hi{उसने} \hi{कमरा} \hi{बुहारा}, \emph{usne kamrā buhārā}.

\begin{cousinweb}
\textbf{\hi{बुहारना}} (\emph{buhārnā}) is a tadbhava: CDIAL files it under Sanskrit \emph{bahukāra}, \textquotedblleft{}effecting much\textquotedblright{}, whose feminine \emph{bahukarī} is \textquotedblleft{}a broom\textquotedblright{}, and gives the Hindi verb as \textquotedblleft{}to sweep\textquotedblright{} (CDIAL 9188).
\end{cousinweb}

\begin{grammarlens}[title={What you've built}]
One more word for describing what you lived and felt.
\end{grammarlens}

\subsection*{Your turn}
\begin{itemize}
\raggedright
  \item \emph{Say it:} \emph{buhārnā}
  \item \emph{Say it:} \emph{buhārnā}, once more
  \item \emph{Say it:} \emph{Dādī har subah āṁgan buhārtī haiṁ.} — Grandmother sweeps the courtyard every morning.
  \item \emph{Recall:} say the Hindi for a spider, then the Hindi for waste paper, then say \emph{buhārnā} again
\end{itemize}

\begin{tcolorbox}[breakable,colback=teal!4,colframe=teal!35!black,title={Before you move on}]
What does \hi{बुहारना} mean? (\textquotedblleft{}To sweep (with a broom)\textquotedblright{}.) Say it once more.
\end{tcolorbox}
\section[\texorpdfstring{\hi{माँजना} (mā̃jnā)}{माँजना (mā̃jnā)}]{\hi{माँजना} (mā̃jnā) — to scour (pots); to scrub clean}

\label{lesson:HI-C519-majna}



\begin{quote}
Before the new one: say the Hindi for waste paper, then the Hindi for to sweep (with a broom).
\end{quote}

\subsection*{You'll want to know: \hi{माँजना}}
\textbf{\hi{माँजना}} — \emph{mā̃jnā} — \textquotedblleft{}to scour (pots); to scrub clean\textquotedblright{}.

\begin{quote}
\hi{खाने} \hi{के} \hi{बाद} \hi{मीरा} \hi{ने} \hi{सब} \hi{बर्तन} \hi{माँजे।}
\end{quote}

\begin{quote}
\emph{Khāne ke bād Mīrā ne sab bartan mā̃je.}
\end{quote}

\begin{quote}
After the meal Mira scoured all the pots.
\end{quote}

Past, with \hi{ने} (he or she …): \hi{उसने} \hi{कड़ाही} \hi{माँजी}, \emph{usne kaṛāhī mā̃jī}.

\begin{cousinweb}
\textbf{\hi{माँजना}} (\emph{mā̃jnā}) is a tadbhava, inherited from Sanskrit \emph{mārjati}, \textquotedblleft{}wipes, cleans\textquotedblright{} (CDIAL 10080). CDIAL gives the Hindi verb as \textquotedblleft{}to rub, clean, polish\textquotedblright{}.
\end{cousinweb}

\begin{grammarlens}[title={What you've built}]
One more word for describing what you lived and felt.
\end{grammarlens}

\subsection*{Your turn}
\begin{itemize}
\raggedright
  \item \emph{Say it:} \emph{mā̃jnā}
  \item \emph{Say it:} \emph{mā̃jnā}, once more
  \item \emph{Say it:} \emph{Khāne ke bād Mīrā ne sab bartan mā̃je.} — After the meal Mira scoured all the pots.
  \item \emph{Recall:} say the Hindi for waste paper, then the Hindi for to sweep (with a broom), then say \emph{mā̃jnā} again
\end{itemize}

\begin{tcolorbox}[breakable,colback=teal!4,colframe=teal!35!black,title={Before you move on}]
What does \hi{माँजना} mean? (\textquotedblleft{}To scour (pots); to scrub clean\textquotedblright{}.) Say it once more.
\end{tcolorbox}
